
\documentclass[letterpaper,12pt]{article}
\usepackage[margin=1in]{geometry}
\usepackage{setspace}
\usepackage[mathlines,edtable]{lineno}
\usepackage{amsmath,amssymb,booktabs,xurl,float}
\usepackage[authoryear,round]{natbib}
\usepackage{microtype}
\usepackage[colorlinks=true,allcolors=blue]{hyperref}
\doublespacing
\renewcommand\linenumberfont{\normalfont\small\sffamily}
\title{Supporting calculations for\\The value of prior information for final dose selection in BOIN trials}
\author{Siyi Chen}
\date{November 2026}
\begin{document}
\maketitle
\linenumbers
This note gives the additional model specifications, comparisons and
computational records for the combined article. The article contains
the method and its proofs. The accompanying data define the evaluated
rules. Fresh searches can return different feasible rules because
optimization is time limited.

\section{Planning information and available gain}

The six-dose study uses target $\varphi=.25$, cohorts of three and at
most 30 participants, starting at dose 1. Its planning support contains
144 curves with positive weight and additional guard curves with zero
planning weight. Location weights are
\[
 q^{\rm central}=(.02,.04,.42,.42,.08,.02),\qquad
 q^{\rm high}=(.01,.02,.04,.08,.20,.65).
\]
The third specification weights the six positions equally.
Within each position, the nominal conditional distribution preserves
the supplied curve weights.

Let $H_{jf}$ denote the conditional distribution for shape family $f$
at MTD position $j$. Only families represented with positive weight at
that position enter the class. The support has 55 such groups.
For a fixed rule $T$, define
\[
 z_{jf}(T)=E_{H_{jf}}\{P_p(T=j)-P_p(b=j)\}.
\]
Allowing a separate mixture of these families at every position gives
the smallest gain
\[
 \min_s\sum_j q_j^s\min_f z_{jf}(T).
\]
The high-location analysis uses one location specification. The joint
analysis uses the convex hull of central, high and equal location
weights. These classes preserve the represented distributions within
each family. They do not include arbitrary new dose--toxicity shapes.

Construction introduces one auxiliary value for each position,
bounded above by every corresponding family gain. A further auxiliary
value is bounded by each weighted location sum. Maximizing the latter
value gives the stated worst-family objective. Risk constraints and
eligibility are the same as in the nominal study. Endpoint
recommendations and one reference outcome per dose remain fixed.

For the upper bound on attainable gain, a nonnegative mixture of the
location vertices and normalized family weights gives an explicit
member of the planning class. A bound on gain under that member also
bounds the smallest gain over the class. The independent calculation
checks this membership, the permitted choices and all nonnegative risk
prices before applying the weak-duality inequality in the article.
The resulting bounds cover deterministic and randomized rules under
the stated restrictions.

\begin{table}[H]
\centering\small\doublespacing
{\internallinenumbers
\caption{Worst-family planning gain and upper bounds, percentage points.
Construction upper is the numerical search bound. Overall upper uses
independent likelihood enclosures and exact rational arithmetic.
Upper bounds are rounded upward.}
\par}
\begin{edtable}{tabular}{lrrr}
\toprule
Planning class & Attained & Construction upper & Overall upper\\
\midrule
High locations, varying shapes & .513 & .520 & .891\\
Varying locations and shapes & .116 & .132 & .261\\
\bottomrule
\end{edtable}
\end{table}

The supplied implementation encloses terminal probabilities with
outward-rounded integer arithmetic. Prior weights are interpreted as
their saved binary values and normalized with exact rational
arithmetic. Signed summation then gives lower and upper gain bounds.
The risk and gain checks are independent of the numerical optimizer.
The construction upper bounds are numerical optimization results.

\section{Bayesian dose--response comparison}

The logistic comparison adapts the posterior-mean final selection of
\citet{wakayama2025}. Their study uses six doses, target .30, 36
participants, three link functions and Hamiltonian Monte Carlo.
The present comparison uses target .25, 30 participants and the logit
link on doses $(7,14,29,51,103,188)$ mg. The model is
\[
 \operatorname{logit}(p_j)=\beta_0+\exp(\beta_1)\log(d_j/51),
 \qquad
 \beta_0\sim N(\mu_0,\sigma_0^2),\quad
 \beta_1\sim N(\mu_1,\sigma_1^2),
\]
with independent priors.

Endpoint beliefs are $P(p_1>\varphi)=.05$ and
$P(p_6\le\varphi)=.21$. Beta marginals express these beliefs. Their
medians are interpolated on the logit scale across log doses.
The model prior is fitted to the .025, .5 and .975 marginal toxicity
quantiles. Calibration depends on the physical grid and the induced
toxicity distributions \citep{neuenschwander2008critical,alger2026calibration}.
A broader prior fixes the log-slope standard deviation at .75 and
refits the other parameters. Table~\ref{tab:logit-priors} summarizes
the two fitted priors. The supplied inputs retain the full precision.

\begin{table}[H]
\centering\small\doublespacing
{\internallinenumbers
\caption{Calibrated logistic priors. Standard deviations are on the
intercept and log-slope scales.}
\label{tab:logit-priors}
\par}
\begin{edtable}{tabular}{lrrrr}
\toprule
Prior & $\mu_0$ & $\sigma_0$ & $\mu_1$ & $\sigma_1$\\
\midrule
Quantile & -1.171 & 1.320 & .102 & .010\\
Broader & -1.270 & 1.201 & .050 & .750\\
\bottomrule
\end{edtable}
\end{table}

Gaussian quadrature at two integration orders checks the posterior
mean toxicities and proposed actions. The calculation uses 181 and
241 nodes per dimension. A separate integral over log slope checks
the fitted marginal quantiles.
The posterior likelihood uses the complete terminal counts.
Path coefficients cancel in posterior normalization.

At each terminal outcome, the proposal selects the eligible dose
whose posterior mean toxicity is closest to target. Endpoint choices,
reference outcomes and no-recommendation status are then retained.
The risk filter chooses either that proposal or BOIN. Central and
high-location filters have their respective nominal planning
objectives. Both logistic priors use the same grid, eligibility,
risk limits and six requested eight-second search rounds as the
finite-prior comparisons.

The three independent evaluation libraries each contain 576 curves:
24 in each of four shape families at every MTD position. Rules are
fixed before these libraries are generated. Table D.1 of the combined
article reports mean PCS and mean gain across the libraries, with the
smallest curve-level gain across all three. These are comparisons of
the specified adaptations under common risk requirements.

\section{Uncertainty from generating evaluation curves}

Terminal-outcome summation gives exact operating characteristics for
a fixed toxicity curve, up to numerical evaluation. The primary
average also varies because its curves are sampled. Each of its four
generator families supplies 40 independent curves at each of six MTD
positions. For paired differences $d_{jfi}$,
\[
 \widehat\Delta=\sum_{jf}(q_j/4)\bar d_{jf},\qquad
 {\rm SE}(\widehat\Delta)=
 \left\{\sum_{jf}(q_j/4)^2s_{jf}^2/n_{jf}\right\}^{1/2}.
\]
Table~\ref{tab:intervals} uses the normal approximation
$\widehat\Delta\pm1.96\,{\rm SE}$. It describes sampling under these
generators. It does not measure uncertainty about a future compound
or about clinical adoption.

\begin{table}[H]
\centering\small\doublespacing
{\internallinenumbers
\caption{Primary mean PCS gains and approximate 95\% generator
intervals, percentage points.}
\label{tab:intervals}
\par}
\begin{edtable}{tabular}{llrr}
\toprule
Prior & Rule & Mean gain & Interval\\
\midrule
Central & Full table & .575 & $(.505,.646)$\\
Central & 80\% table & .522 & $(.467,.577)$\\
Central & 90\% table & .578 & $(.517,.638)$\\
High & Full table & .522 & $(.491,.553)$\\
High & 80\% table & .396 & $(.366,.427)$\\
High & 90\% table & .461 & $(.430,.492)$\\
\bottomrule
\end{edtable}
\end{table}

\section{Comparison of risk bounds}

The ordered-likelihood and Continuous Simulation Extension (CSE)
calculations start from the same 1,716 cells and target a two-point
risk bound. They use the same fixed central and high rules. Interior
PCS-loss and both signs of each tail contrast are checked. The CSE
implementation uses exact terminal expectations at each knowledge
point \citep{yang2024cse}.

Write $L$ for the likelihood ratio from knowledge point $p^0$ to $p$.
For any event $A$ and $r>1$, the tilt bound gives
\[
 P_p(A)\le
 P_{p^0}(A)^{1-1/r}
 \{E_{p^0}(L^r)\}^{1/r}.
\]
Positive and negative events in a signed contrast are bounded
separately. The negative-event lower bound uses both its complement
and the reverse likelihood ratio. An alternative applies the tilt
bound to the contrast shifted and scaled between zero and one.
The calculation retains the smallest valid upper bound.

The independent-sequence embedding supplies each dose with its
largest reachable enrollment. A second embedding accounts for
bounded stopping. For a cell with coordinate bounds
$(\ell_j,u_j)$, define
\[
 \bar\psi_j=\max_{t\in[\ell_j,u_j]}
 \log\{t^r(p_j^0)^{1-r}+
          (1-t)^r(1-p_j^0)^{1-r}\}.
\]
The quantity before taking the logarithm is convex, so the coordinate
maximum occurs at an endpoint. A nonnegative supermartingale and
bounded stopping give
\[
 E_{p^0}(L_\tau^r)\le
 \exp\left\{\max_{n\in\mathcal N}\sum_j n_j\bar\psi_j\right\},
\]
where $\mathcal N$ contains the reachable final enrollment vectors.
The reverse calculation exchanges the two probability measures.

All methods bisect unresolved cells along their widest coordinate.
Children retain the smaller of their new bound and their parent's
bound. The original comparison gives each method 60 seconds after
initialization and at most 200,000 leaves. Its timings and refinement
histories are supplied. Replaying to the recorded final refinement
counts checks the reported bounds without making results depend on
the speed of a later computer. A separate command repeats the
resource-limited benchmark on the current computer. Independent
moment checks compare the stopped-trial bound with direct terminal
likelihood-ratio sums.

CSE results are numerical tilt calculations. The completed
ordered-likelihood calculation also receives the independent integer
risk check used for the final selection rules. Table 6 of the article
reports the original comparison.

\section{Other trial designs and computational records}

The supporting designs vary maximum enrollment, grid size and cohort
size. They use target .25 and starting dose 1. Their planning curves
have the form
\[
 p_i=\operatorname{logit}^{-1}
 \{\operatorname{logit}(\varphi)+s(i-j+t)\}.
\]
The six slopes are $.3,.5,.8,1.2,1.8,2.6$. Interior offsets are
$\pm.08$ and $\pm.22$. Endpoint offsets point inward with magnitudes
$.04,.08,.15,.22$. The 24 curves at each position have equal
conditional weight. Six-dose location weights match the main study.
The four-dose weights are $(.04,.46,.46,.04)$ and
$(.03,.08,.24,.65)$. The eight-dose weights are
\[
\begin{split}
 q^{\rm central}&=(.02,.03,.10,.35,.35,.10,.03,.02),\\
 q^{\rm high}&=(.01,.01,.02,.03,.05,.08,.15,.65).
\end{split}
\]

\begin{table}[H]
\centering\small\doublespacing
{\internallinenumbers
\caption{Supporting designs on fixed descriptive evaluation libraries.
BOIN PCS is a percentage. Signed and Separate are changes in percentage
points. Worst is the smallest curve-level change for the signed rule.}
\par}
\begin{edtable}{tabular}{llrrrr}
\toprule
Design & Prior & BOIN PCS & Signed & Separate & Worst\\
\midrule
6 doses, 36 participants & Central & 35.159 & .660 & .545 & -1.701\\
6 doses, 36 participants & High & 29.877 & .305 & .294 & -1.263\\
8 doses, 24 participants & Central & 26.162 & .363 & .367 & -1.498\\
8 doses, 24 participants & High & 13.404 & .055 & .053 & -1.370\\
4 doses, cohorts of 2 & Central & 36.547 & -.004 & .013 & -1.047\\
4 doses, cohorts of 2 & High & 35.032 & .697 & .494 & -1.622\\
\bottomrule
\end{edtable}
\end{table}

Signed rules use both signs of each risk contrast within ordered
cells. Separate rules bound adverse contributions state by state.
The supporting calculations rebuild each catalog, check every saved
rule, and recompute its operating characteristics and planning gains.
The eight-dose partition starts from five coordinate intervals.
The other supporting partitions start from eight. Each final
partition has 1,716 cells. The broader search can retain a feasible
comparator when the time-limited search does not find a better rule.

The original nominal workflow takes 833.3 seconds, including
initialization, enumeration, prior integration, optimization,
verification and decision-table export. Its largest child process
uses 4,588 MiB. Given those frozen inputs, the endpoint, shorter-table,
CSE and validation stage takes 885.1 seconds with peak child memory
of 3,647 MiB. The added robust-prior and clinical analyses have their
own timing records. The application stage takes 33.3 seconds from
its prepared inputs. These stages have different starting inputs
and are not interchangeable timing estimates.

The original measurements are retained as records of those
workflows. The current R/Rcpp replay records its own phase times,
software versions and hardware. It recomputes operating
characteristics and verifies the frozen rules and gain bounds.
Fresh fitting is a separate operation with explicit search limits.
The saved tables define the article's results.

The package supplies the complete decision tables, planning
distributions, evaluation curves, risk and gain inputs, and a mapping
from every article table to its calculation. All analysis and plotting
programs use R or Rcpp. Final trial use requires only the observed
counts and the prespecified table.

\begin{thebibliography}{9}
\bibitem[Alger et al.(2026)]{alger2026calibration}
Alger E, Lee SM, Cheung YKK, Yap C.
Calibration of priors for Bayesian model-based dose-finding trial
designs with joint outcomes. \emph{Statistics in Medicine},
45(23--24):e70746, 2026. doi:10.1002/sim.70746.

\bibitem[Neuenschwander et al.(2008)]{neuenschwander2008critical}
Neuenschwander B, Branson M, Gsponer T.
Critical aspects of the Bayesian approach to phase I cancer trials.
\emph{Statistics in Medicine}, 27(13):2420--2439, 2008.
doi:10.1002/sim.3230.

\bibitem[Wakayama et al.(2025)]{wakayama2025}
Wakayama R, Momozaki T, Ando S.
Improving maximum tolerated dose selection in model-assisted designs
for phase I trials through Bayesian dose-response model.
\emph{Pharmaceutical Statistics}, 24(6):e70046, 2025.
doi:10.1002/pst.70046.

\bibitem[Yang et al.(2024)]{yang2024cse}
Yang J, Thompson TB, Sklar M.
Guarantees for comprehensive simulation assessment of statistical
methods. \emph{arXiv preprint}, arXiv:2212.10042v3, 2024.
doi:10.48550/arXiv.2212.10042.
\end{thebibliography}
\end{document}
